\scanpage{91}
% archive foliation: 90
disjoints. Donc la donnée d'un homomorphisme de groupes
$G\quad\uline{G}_{m\,S}$
équivaut à la donnée d'une décomposition de $S$ comme réunion d'ouverts
$S_m\ (m\quad M)$ deux à deux disjoints, i.e. à la donnée d'une application
localement constante de $S$ dans $M$. Cela établit 1.3. donc 1.2.

\uline{Corollaire} 1.4. Un groupe diagonalisable est réflexif, donc aussi un
groupe localement diagonalisable. Si $M,N$ sont deux groupes commutatifs
ordinaires, \deleted{\unsure{alors}} l'homomorphisme naturel
\[
\operatorname{Hom}_{S\text{-gr}}(M_S,N_S)
\qquad
\operatorname{Hom}_{S\text{-gr}}(D(N_S),D(M_S))
\]
est bijectif.

Notons maintenant que si $M$ est un groupe commutatif \uline{de type fini},
\deleted{\unsure{il est}} donc défini par un nombre fini de générateurs et un nombre fini de
relations, alors pour tout
\deleted{\unsure{préschéma en groupes $G$ sur $S$, on a une}}
\inserted{groupe ordinaire $N$,}
l'application naturelle

On en déduit un isomorphisme de foncteurs sur $S$:
\[
\uline{\operatorname{Hom}}_{S\text{-gr}}(M_S,N_S)
\qquad
\uline{\operatorname{Hom}}_{S\text{-gr}}(D(N_S),D(M_S))
\]
D'autre part \unsure{un|on} définit de façon évidente un homomorphisme
\[
\operatorname{Hom}_{\mathrm{gr}}(M,N)_S
\qquad
\uline{\operatorname{Hom}}_{S\text{-gr}}(M_S,N_S),
\]
qui est toujours \deleted{\unsure{inj|un}} un monomorphisme, et qui est un isomorphisme dès
que $M$ est de type fini. On trouve ainsi:

\uline{Corollaire} 1.5. \deleted{\unsure{Supposons $M$}}
Soient $M,N$ deux groupes commutatifs
ordinaires, $M$ de type fini, alors \deleted{\unsure{on a}}
$\uline{\operatorname{Hom}}_{S\text{-gr}}(D(N_S),D(M_S))$ est
représentable, et on a un isomorphisme
\[
\uline{\operatorname{Hom}}_{\mathrm{gr}}(M,N)_S
\qquad
\uline{\operatorname{Hom}}_{S\text{-gr}}(D(N_S),D(M_S))
\]
On en conclut plus généralement que si $G,H$ sont localement diagonalisables,
$H$ étant de type fini, alors
$\uline{\operatorname{Hom}}_{S\text{-gr}}(G,H)$ est représentable.
